\documentclass{beamer}
\usetheme{Madrid}
\usepackage{graphicx}
\usepackage{amsmath}
\usepackage{booktabs}
\usepackage{xcolor}
\title{Attention Is All You Need}
\author{Ashish VaswaniEqual contribution. Listing order is random. Jakob proposed replacing RNNs with self-attention and started the effort to evaluate this idea. Ashish, with Illia, designed and implemented the first Transformer models and has been crucially involved in every aspect of this work. Noam proposed scaled dot-product attention, multi-head attention and the parameter-free position representation and became the other person involved in nearly every detail. Niki designed, implemented, tuned and evaluated countless model variants in our original codebase and tensor2tensor. Llion also experimented with novel model variants, was responsible for our initial codebase, and efficient inference and visualizations. Lukasz and Aidan spent countless long days designing various parts of and implementing tensor2tensor, replacing our earlier codebase, greatly improving results and massively accelerating our research.}
\definecolor{c1}{HTML}{FFFFFF}\n\definecolor{c2}{HTML}{333333}\n\definecolor{c3}{HTML}{2E7D32}\n\begin{document}
\frame{\titlepage}
\begingroup
\setbeamercolor{background canvas}{bg=c1}\n\setbeamercolor{normal text}{fg=c2}\n\begin{frame}[]{Introduction to Sequence Modeling}\n\\begin{itemize}
\\item \\alert{Recurrent Neural Networks (RNNs) are state-of-the-art for sequence modeling.}\\n\\item \\alert{Long Short-Term Memory (LSTM) and Gated Recurrent Units (GRUs) enhance RNNs.}\\n\\item Sequential nature limits parallelization, especially with long sequences.\\n\\end{itemize}
\\end{frame}

\endgroup

\begingroup
\setbeamercolor{background canvas}{bg=c1}\n\setbeamercolor{normal text}{fg=c2}\n\begin{frame}[]{Challenges with Recurrent Models}\n\\begin{itemize}
\\item \\alert{RNNs compute hidden states sequentially: h\_t = f(h\_t-1, input\_t).}\\n\\item \\alert{This leads to linear time complexity with respect to input/output length.}\\n\\item Memory constraints limit batching across examples during training.\\n\\end{itemize}
\\end{frame}

\endgroup

\begingroup
\setbeamercolor{background canvas}{bg=c1}\n\setbeamercolor{normal text}{fg=c2}\n\begin{frame}[]{The Role of Attention Mechanisms}\n\\begin{itemize}
\\item \\alert{Attention mechanisms model dependencies regardless of distance.}\\n\\item \\alert{Traditionally used with recurrent networks for improved performance.}\\n\\item Emergence of attention-based models has transformed machine translation.\\n\\end{itemize}
\\end{frame}

\endgroup

\begingroup
\setbeamercolor{background canvas}{bg=c1}\n\setbeamercolor{normal text}{fg=c2}\n\begin{frame}[]{Introducing the Transformer Model}\n\\begin{itemize}
\\item \\alert{The Transformer architecture relies entirely on attention mechanisms.}\\n\\item \\alert{Eliminates recurrence, enabling significant parallelization.}\\n\\item Achieves state-of-the-art translation quality in reduced training time.\\n\\end{itemize}
\\end{frame}

\endgroup

\begingroup
\setbeamercolor{background canvas}{bg=c4}\n\setbeamercolor{normal text}{fg=c5}\n\begin{frame}[plain]{Understanding RNNs Analogy}\n\\begin{itemize}
\\item \\alert{Imagine a teacher (RNN) helping students (data) learn a story (sequence) one sentence at a time.}\\n\\item The teacher can only focus on the last sentence heard, making it hard to remember earlier parts.\\n\\end{itemize}
\\vspace{3mm}
\\begin{center}
\\begin{beamercolorbox}[rounded=true,shadow=true]{title}
Use friendly graphics of a teacher and students.
\\end{beamercolorbox}
\\end{center}
\\end{frame}

\endgroup

\begingroup
\setbeamercolor{background canvas}{bg=c7}\n\setbeamercolor{normal text}{fg=c8}\n\begin{frame}[]{Questions for Discussion}\n\\begin{itemize}
\\item How do attention mechanisms improve upon traditional RNNs?\\n\\item What are the implications of eliminating recurrence in sequence models?\\n\\end{itemize}
\\end{frame}

\endgroup

\begingroup
\setbeamercolor{background canvas}{bg=c1}\n\setbeamercolor{normal text}{fg=c2}\n\begin{frame}[]{Background of Transformer Models}\n\\begin{itemize}
\\item \\alert{Goal: Reduce sequential computation in neural networks.}\\n\\item \\alert{Models: Extended Neural GPU, ByteNet, ConvS2S utilize CNNs for parallel computation.}\\n\\item \\alert{Operation growth: Linear (ConvS2S) vs. logarithmic (ByteNet) with distance.}\\n\\item Transformers reduce operations to constant time, averaging attention-weighted positions.\\n\\item Self-attention relates different positions in a sequence for enhanced representation.\\n\\end{itemize}
\\end{frame}

\endgroup

\begingroup
\setbeamercolor{background canvas}{bg=c4}\n\setbeamercolor{normal text}{fg=c5}\n\begin{frame}[plain]{analogy: Understanding Self-Attention}\n\\begin{itemize}
\\item Imagine reading a book and using a highlighter to mark important sentences.\\n\\item Self-attention is like that highlighter, helping the model focus on relevant parts of the text.\\n\\item Just as you can highlight multiple sentences at once, self-attention connects different words in a sentence.\\n\\end{itemize}
\\end{frame}

\endgroup

\begingroup
\setbeamercolor{background canvas}{bg=c1}\n\setbeamercolor{normal text}{fg=c2}\n\begin{frame}[]{Applications of Self-Attention}\n\\begin{itemize}
\\item \\alert{Used in tasks: Reading comprehension, summarization, entailment.}\\n\\item \\alert{End-to-end memory networks utilize recurrent attention mechanisms.}\\n\\item Transformers: The first models using only self-attention for input-output representation.\\n\\end{itemize}
\\end{frame}

\endgroup

\begingroup
\setbeamercolor{background canvas}{bg=c1}\n\setbeamercolor{normal text}{fg=c2}\n\begin{frame}[]{Efficiency of Attention Mechanisms}\n\\begin{itemize}
\\item \\alert{Attention layers connect numerous positions at low computation cost.}\\n\\item \\alert{Question: Can recurrence be replaced with attention?}\\n\\item Pure attention models achieve state-of-the-art results in machine translation.\\n\\end{itemize}
\\end{frame}

\endgroup

\begingroup
\setbeamercolor{background canvas}{bg=c4}\n\setbeamercolor{normal text}{fg=c5}\n\begin{frame}[plain]{analogy: Visualizing Attention Efficiency}\n\\begin{itemize}
\\item Think of a librarian who can instantly find any book in a library.\\n\\item Attention mechanisms work like that librarian, quickly connecting relevant information without searching through every book.\\n\\item This efficiency allows models to process information faster and more effectively.\\n\\end{itemize}
\\end{frame}

\endgroup

\begingroup
\setbeamercolor{background canvas}{bg=c7}\n\setbeamercolor{normal text}{fg=c8}\n\begin{frame}[]{Questions for Discussion}\n\\begin{itemize}
\\item How does self-attention improve upon traditional RNNs?\\n\\item What are the implications of reducing sequential computation?\\n\\end{itemize}
\\end{frame}

\endgroup

\begingroup
\setbeamercolor{background canvas}{bg=c1}\n\setbeamercolor{normal text}{fg=c2}\n\begin{frame}[]{Model Architecture Overview}\n\\begin{itemize}
\\item \\alert{Encoder-decoder structure common in sequence-to-sequence models}\\n\\item \\alert{Encoder computes representation of input sequence}\\n\\item Decoder autoregressively produces output sequence\\n\\item Utilizes attention layers and position-wise feed-forward layers\\n\\end{itemize}
\\end{frame}

\endgroup

\begingroup
\setbeamercolor{background canvas}{bg=c1}\n\setbeamercolor{normal text}{fg=c2}\n\begin{frame}[]{Encoder and Decoder Components}\n\\begin{itemize}
\\item \\alert{Encoder: Stack of N=6 identical layers}\\n\\item \\alert{Each layer includes: Multi-head self-attention, Position-wise feed-forward network}\\n\\item Decoder: Similar stack with an additional attention layer\\n\\item Residual connections and layer normalization are applied\\n\\end{itemize}
\\end{frame}

\endgroup

\begingroup
\setbeamercolor{background canvas}{bg=c1}\n\setbeamercolor{normal text}{fg=c2}\n\begin{frame}[]{Attention Mechanism}\n\\begin{itemize}
\\item \\alert{Maps queries and key-value pairs to outputs}\\n\\item \\alert{Output is a weighted sum of values based on compatibility with keys}\\n\\item Scaled Dot-Product Attention enhances efficiency\\n\\item Multi-head attention allows parallel processing of information\\n\\end{itemize}
\\end{frame}

\endgroup

\begingroup
\setbeamercolor{background canvas}{bg=c4}\n\setbeamercolor{normal text}{fg=c5}\n\begin{frame}[plain]{Attention Mechanism Analogy}\n\\begin{itemize}
\\item \\alert{Imagine a librarian (query) looking for books (keys) on a shelf (values).}\\n\\item The librarian checks each book's title (compatibility) to find the right one.\\n\\item The more relevant the title, the more attention the librarian gives to that book.\\n\\end{itemize}
\\vspace{3mm}
\\begin{center}
\\begin{beamercolorbox}[rounded=true,shadow=true]{title}
Use friendly graphics of a librarian and books
\\end{beamercolorbox}
\\end{center}
\\end{frame}

\endgroup

\begingroup
\setbeamercolor{background canvas}{bg=c1}\n\setbeamercolor{normal text}{fg=c2}\n\begin{frame}[]{Positional Encoding}\n\\begin{itemize}
\\item \\alert{Positional encodings added to input embeddings to retain order information}\\n\\item Sine and cosine functions are used for encoding positions\\n\\item Enables the model to learn relative positions effectively\\n\\end{itemize}
\\end{frame}

\endgroup

\begingroup
\setbeamercolor{background canvas}{bg=c4}\n\setbeamercolor{normal text}{fg=c5}\n\begin{frame}[plain]{Positional Encoding Analogy}\n\\begin{itemize}
\\item \\alert{Think of a train with cars (values) that need to stay in order.}\\n\\item Each car has a number (position) that tells you where it belongs.\\n\\item Without numbers, the train would be a jumbled mess!\\n\\end{itemize}
\\vspace{3mm}
\\begin{center}
\\begin{beamercolorbox}[rounded=true,shadow=true]{title}
Use cheerful graphics of a train with numbered cars
\\end{beamercolorbox}
\\end{center}
\\end{frame}

\endgroup

\begingroup
\setbeamercolor{background canvas}{bg=c7}\n\setbeamercolor{normal text}{fg=c8}\n\begin{frame}[]{Questions for Discussion}\n\\begin{itemize}
\\item How does the attention mechanism improve translation tasks?\\n\\item What are the advantages of using multi-head attention?\\n\\item Why is positional encoding necessary in this model?\\n\\end{itemize}
\\end{frame}

\endgroup

\begingroup
\setbeamercolor{background canvas}{bg=c1}\n\setbeamercolor{normal text}{fg=c2}\n\begin{frame}[]{Why Self-Attention?}\n\\begin{itemize}
\\item \\alert{Mapping variable-length sequences: x to y}\\n\\item \\alert{Comparison with recurrent and convolutional layers}\\n\\item \\alert{Three key factors for self-attention:}\\n\\item 1. Computational complexity per layer\\n\\item 2. Parallelizable computation\\n\\item 3. Path length for long-range dependencies\\n\\end{itemize}
\\end{frame}

\endgroup

\begingroup
\setbeamercolor{background canvas}{bg=c1}\n\setbeamercolor{normal text}{fg=c2}\n\begin{frame}[]{Self-Attention vs. Other Layers}\n\\begin{itemize}
\\item \\alert{Self-Attention: O(n d), O(1), O(1)}\\n\\item \\alert{Recurrent: O(n d), O(n), O(n)}\\n\\item \\alert{Convolutional: O(k n d), O(1), O(log\_k(n))}\\n\\item Efficiency in cross-positional communications\\n\\end{itemize}
\\end{frame}

\endgroup

\begingroup
\setbeamercolor{background canvas}{bg=c1}\n\setbeamercolor{normal text}{fg=c2}\n\begin{frame}[]{Path Length and Receptive Field}\n\\begin{itemize}
\\item \\alert{Receptive field size influences output representations}\\n\\item \\alert{Self-Attention: full receptive field = n}\\n\\item \\alert{Recurrent: full receptive field = n}\\n\\item Convolutional: limited receptive field = k\\n\\item Shorter paths facilitate learning of long-range dependencies\\n\\end{itemize}
\\end{frame}

\endgroup

\begingroup
\setbeamercolor{background canvas}{bg=c4}\n\setbeamercolor{normal text}{fg=c5}\n\begin{frame}[plain]{Path Length Analogy}\n\\begin{itemize}
\\item \\alert{Think of a race: shorter paths help runners reach the finish line faster.}\\n\\item Self-attention allows all runners (data points) to connect directly, speeding up learning.\\n\\item Recurrent layers are like a relay race, passing the baton one at a time.\\n\\item Convolutional layers are like a team of runners only communicating in small groups.\\n\\end{itemize}
\\vspace{3mm}
\\begin{center}
\\begin{beamercolorbox}[rounded=true,shadow=true]{title}
Use race graphics or running icons
\\end{beamercolorbox}
\\end{center}
\\end{frame}

\endgroup

\begingroup
\setbeamercolor{background canvas}{bg=c4}\n\setbeamercolor{normal text}{fg=c5}\n\begin{frame}[plain]{Self-Attention Analogy}\n\\begin{itemize}
\\item \\alert{Imagine a classroom where every student can talk to every other student at the same time.}\\n\\item This is like self-attention: everyone shares ideas directly.\\n\\item In contrast, a teacher (recurrent layer) can only talk to one student at a time.\\n\\item A group project (convolutional layer) only allows a few students to communicate.\\n\\end{itemize}
\\vspace{3mm}
\\begin{center}
\\begin{beamercolorbox}[rounded=true,shadow=true]{title}
Use friendly classroom graphics or lightbulb icon
\\end{beamercolorbox}
\\end{center}
\\end{frame}

\endgroup

\begingroup
\setbeamercolor{background canvas}{bg=c7}\n\setbeamercolor{normal text}{fg=c8}\n\begin{frame}[]{Questions for Discussion}\n\\begin{itemize}
\\item How does self-attention improve learning of long-range dependencies?\\n\\item What are the trade-offs between self-attention and recurrent layers?\\n\\item Can you think of scenarios where convolutional layers might be preferable?\\n\\end{itemize}
\\end{frame}

\endgroup

\begingroup
\setbeamercolor{background canvas}{bg=c1}\n\setbeamercolor{normal text}{fg=c2}\n\begin{frame}[]{Training Regime Overview}\n\\begin{itemize}
\\item \\alert{Training on WMT 2014 English-German dataset: \textasciitilde{}4.5M sentence pairs}\\n\\item \\alert{WMT 2014 English-French dataset: 36M sentences}\\n\\item Byte-pair encoding for tokenization: 37,000 tokens (English-German), 32,000 tokens (English-French)\\n\\item Batching by approximate sequence length: \textasciitilde{}25,000 source \& target tokens\\n\\end{itemize}
\\end{frame}

\endgroup

\begingroup
\setbeamercolor{background canvas}{bg=c1}\n\setbeamercolor{normal text}{fg=c2}\n\begin{frame}[]{Training Details}\n\\begin{itemize}
\\item \\alert{Hardware: 8 NVIDIA P100 GPUs}\\n\\item \\alert{Base model training: 100,000 steps (\textasciitilde{}12 hours), step time: 0.4 seconds}\\n\\item Big model training: 300,000 steps (\textasciitilde{}3.5 days), step time: 1.0 seconds\\n\\item Optimizer: Adam with parameters (1=0.9, 2=0.98, =10\textasciicircum{}-9)\\n\\end{itemize}
\\end{frame}

\endgroup

\begingroup
\setbeamercolor{background canvas}{bg=c1}\n\setbeamercolor{normal text}{fg=c2}\n\begin{frame}[]{Learning Rate Schedule}\n\\begin{itemize}
\\item \\alert{Learning rate formula: lrate = step\_num\textasciicircum{}-0.5 * (step\_num warmup\_steps\textasciicircum{}-1.5)}\\n\\item \\alert{Warmup steps: 4,000}\\n\\item Initial increase followed by decrease proportional to inverse square root of step number\\n\\end{itemize}
\\end{frame}

\endgroup

\begingroup
\setbeamercolor{background canvas}{bg=c4}\n\setbeamercolor{normal text}{fg=c5}\n\begin{frame}[plain]{Learning Rate Analogy}\n\\begin{itemize}
\\item Think of learning rate like a car's speed: too fast and you might crash (overfitting), too slow and you won't reach your destination (underfitting).\\n\\item \\alert{Warmup steps are like gradually pressing the gas pedal to avoid sudden acceleration.}\\n\\end{itemize}
\\end{frame}

\endgroup

\begingroup
\setbeamercolor{background canvas}{bg=c1}\n\setbeamercolor{normal text}{fg=c2}\n\begin{frame}[]{Regularization Techniques}\n\\begin{itemize}
\\item \\alert{Dropout applied to output of each sub-layer and embeddings}\\n\\item \\alert{Base model dropout rate: P\_drop=0.1}\\n\\item Attention dropout: dropping attention weights during training\\n\\item Label smoothing: value \_ls=0.1 to improve accuracy and BLEU score\\n\\end{itemize}
\\end{frame}

\endgroup

\begingroup
\setbeamercolor{background canvas}{bg=c7}\n\setbeamercolor{normal text}{fg=c8}\n\begin{frame}[]{Questions for Discussion}\n\\begin{itemize}
\\item How does dropout help in training models?\\n\\item What impact does learning rate scheduling have on model performance?\\n\\item Why is label smoothing beneficial despite increasing perplexity?\\n\\end{itemize}
\\end{frame}

\endgroup

\begingroup
\setbeamercolor{background canvas}{bg=c1}\n\setbeamercolor{normal text}{fg=c2}\n\begin{frame}[]{Results Overview}\n\\begin{itemize}
\\item \\alert{Comparison of translation models on EN-DE and EN-FR tasks}\\n\\item \\alert{Metrics: cBLEU and Training Cost (FLOPs)}\\n\\item Transformer models outperform previous state-of-the-art\\n\\end{itemize}
\\end{frame}

\endgroup

\begingroup
\setbeamercolor{background canvas}{bg=c1}\n\setbeamercolor{normal text}{fg=c2}\n\begin{frame}[]{Model Performance Comparison}\n\\begin{itemize}
\\item \\alert{ByteNet: cBLEU = 23.75 (EN-DE)}\\n\\item \\alert{Deep-Att + PosUnk: cBLEU = 39.2 (EN-FR), Training Cost = 1.01e20 FLOPs}\\n\\item \\alert{GNMT + RL: cBLEU = 24.6 (EN-DE), 39.92 (EN-FR)}\\n\\item ConvS2S: cBLEU = 25.16 (EN-DE), 40.46 (EN-FR)\\n\\item Transformer (base model): cBLEU = 27.3 (EN-DE), Training Cost = 3.31e18 FLOPs\\n\\end{itemize}
\\end{frame}

\endgroup

\begingroup
\setbeamercolor{background canvas}{bg=c1}\n\setbeamercolor{normal text}{fg=c2}\n\begin{frame}[]{Transformer Model Achievements}\n\\begin{itemize}
\\item \\alert{Big Transformer model achieves BLEU score of 28.4 on EN-DE}\\n\\item \\alert{Outperforms best previously reported models by over 2.0 BLEU}\\n\\item Training duration: 3.5 days on 8 P100 GPUs\\n\\item Big model achieves BLEU score of 41.0 on EN-FR, less than 1/4 training cost of previous state-of-the-art\\n\\end{itemize}
\\end{frame}

\endgroup

\begingroup
\setbeamercolor{background canvas}{bg=c1}\n\setbeamercolor{normal text}{fg=c2}\n\begin{frame}[]{Hyperparameter Tuning Insights}\n\\begin{itemize}
\\item \\alert{Dropout rate: P\_drop = 0.1 for big model, 0.3 for base model}\\n\\item \\alert{Beam search with size 4 and length penalty = 0.6}\\n\\item Maximum output length set to input length + 50\\n\\end{itemize}
\\end{frame}

\endgroup

\begingroup
\setbeamercolor{background canvas}{bg=c4}\n\setbeamercolor{normal text}{fg=c5}\n\begin{frame}[plain]{Analogy: Understanding the Transformer}\n\\begin{itemize}
\\item Think of the Transformer as a smart translator that learns from lots of books.\\n\\item It understands words and their meanings by looking at the context around them.\\n\\item Just like how we learn languages by reading and listening, the Transformer learns from data!\\n\\end{itemize}
\\end{frame}

\endgroup

\begingroup
\setbeamercolor{background canvas}{bg=c7}\n\setbeamercolor{normal text}{fg=c8}\n\begin{frame}[]{Q\&A: Discussion on Model Performance}\n\\begin{itemize}
\\item What factors contribute to the performance of the Transformer model?\\n\\item How does the training cost impact the choice of model in practice?\\n\\item What might be the implications of achieving state-of-the-art BLEU scores?\\n\\end{itemize}
\\end{frame}

\endgroup

\begingroup
\setbeamercolor{background canvas}{bg=c1}\n\setbeamercolor{normal text}{fg=c2}\n\begin{frame}[]{Conclusion of the Transformer Model}\n\\begin{itemize}
\\item \\alert{Introduced the Transformer: a sequence transduction model based entirely on attention.}\\n\\item \\alert{Replaced recurrent layers with multi-headed self-attention.}\\n\\item \\alert{Achieved faster training for translation tasks compared to recurrent or convolutional architectures.}\\n\\item Set new state-of-the-art results on WMT 2014 English-to-German and English-to-French translation tasks.\\n\\item Best model outperformed all previously reported ensembles.\\n\\end{itemize}
\\begin{center}
\\includegraphics[width=0.9\\textwidth]{/Users/dhruv/Documents/Agents/arXiv-1706.03762v7/vis/making_more_difficult5_new.pdf}\\n\\end{center}
\\begin{center}
\\includegraphics[width=0.9\\textwidth]{/Users/dhruv/Documents/Agents/arXiv-1706.03762v7/vis/anaphora_resolution_new.pdf}\\n\\end{center}
\\begin{center}
\\includegraphics[width=0.9\\textwidth]{/Users/dhruv/Documents/Agents/arXiv-1706.03762v7/vis/anaphora_resolution2_new.pdf}\\n\\end{center}
\\begin{center}
\\includegraphics[width=0.9\\textwidth]{/Users/dhruv/Documents/Agents/arXiv-1706.03762v7/vis/attending_to_head_new.pdf}\\n\\end{center}
\\begin{center}
\\includegraphics[width=0.9\\textwidth]{/Users/dhruv/Documents/Agents/arXiv-1706.03762v7/vis/attending_to_head2_new.pdf}\\n\\end{center}
\\end{frame}

\endgroup

\begingroup
\setbeamercolor{background canvas}{bg=c4}\n\setbeamercolor{normal text}{fg=c5}\n\begin{frame}[plain]{Analogy: Understanding the Transformer Model}\n\\begin{itemize}
\\item Think of the Transformer like a group of friends sharing stories.\\n\\item Each friend listens to the others and remembers important parts.\\n\\item Instead of taking turns (like older models), they can all talk at once!\\n\\item This helps them share stories faster and more clearly.\\n\\end{itemize}
\\end{frame}

\endgroup

\begingroup
\setbeamercolor{background canvas}{bg=c1}\n\setbeamercolor{normal text}{fg=c2}\n\begin{frame}[]{Future Directions}\n\\begin{itemize}
\\item \\alert{Broader applicability of models: experiments on English constituency parsing.}\\n\\item \\alert{Plans to extend Transformer to other input/output modalities (e.g., images, audio, video).}\\n\\item Research goals include making generation less sequential and investigating local attention mechanisms.\\n\\end{itemize}
\\end{frame}

\endgroup

\begingroup
\setbeamercolor{background canvas}{bg=c7}\n\setbeamercolor{normal text}{fg=c8}\n\begin{frame}[]{Questions for Discussion}\n\\begin{itemize}
\\item What implications does the Transformer model have for future research?\\n\\item How might attention mechanisms be applied to non-text data?\\n\\end{itemize}
\\end{frame}

\endgroup

\end{document}
